% page 623 du fac-simile (image p-622 du scan)
% bloc 167.6 x 250.6 mm ; marge gauche 34.4 mm
\pgc{12.700mm}{0.000mm}{
\l{\cel{53}615}
\l{}
\l{vestono. Vesto virala, kelke plu kurta kam vestio.- F.}
\l{}
\l{vetar. (trans.) Refuzar sua sanciono di (ulo). EF.}
\l{}
\l{vetero.\cel{2}(meteorol.) La stando di atmosfero. - DE.}
\l{}
\l{veterano. l. (olim, en Roma). Soldato qua obtenabis sua konjedo}
\l{\cel{3}definitiva, pos armeanesir dum multa tempo. - II. (\sou{nun)}.}
\l{\cel{3}Soldato qua esas armeano de multa tempo. - DEFIRS.}
\l{}
\l{veterinaro.\cel{2}Mediko qua kuracas la morbi di la kavali, hundi,}
\l{\cel{3}e c. - DEFIRS.}
\l{}
\l{\sou{vetivero}. (\sou{bot.})\cel{3}Planto de India, uzata por facar parfumi. -}
\l{\cel{3}- L.\cel{1}\sou{andropogon vetiveris}. DEFS.}
\l{}
\l{\sou{vetrino}. (\sou{tekn.}) I. La parto di la butiko qua separesas de la}
\l{\cel{3}strado per nur granda panelo de vitro. - II. Armoro, tablo}
\l{\cel{3}sur qua esas instalita kesto ek vitro, en qua onu expozas la}
\l{\cel{3}objekti skarsa o frajila. - FI.}
\l{}
\l{\sou{veturo}. Vehilo por transporti, ek framo muntita sur roti, quan}
\l{\cel{3}tranas kavali, muli, e c. od homo. - FI.}
\l{}
\l{\sou{vexar}. (\sou{trans}.,\cel{1}\sou{pri},\cel{1}\sou{pro)}.\cel{2}Igar (ulu) deskontenta, despitanta}
\l{\cel{3}e pasable shamoza ed interne iracoza, kande ulu amuzas su per}
\l{\cel{3}kontre-agar lu pri mikraji. - DE.}
\l{}
\l{\sou{veziko.}\cel{2}I. (\sou{anat.})\cel{2}Quaza sako, posho, en qua urino esas akumu-}
\l{\cel{3}le, til kande la parieti kontraktas por ekpulsar lu. - II.}
\l{\cel{3}Saketo en la abdomino di la maxim multa fishi, qua kontenas}
\l{\cel{3}mixuro de azoto, oxigeno e karbo-bioxo, ed igas la fisho plu}
\l{\cel{3}lejera. - III.Ta sako, tirita ek la korpo di la animalo, siki-}
\l{\cel{3}gita ed inflita da la aero quan onu ensuflis. - EFIL.}
\l{}
\l{\sou{vezikatorio.}\cel{1}Substanco qua produktas veziketi sur la pelo. -}
\l{\cel{3}DEFIS.}
\l{}
\l{vi.\cel{2}(\sou{gram.}) Pluralo di \textquotedbl{}vu\textquotedbl{}.}
\l{}
\l{viadukto. (\sou{tekn.)}\cel{2}Ponto alta, kun,multa-kaze,plura rangi de}
\l{\cel{3}arkadi, konstruktita super fluvio, rivero, valeto, e c. -}
\l{\cel{3}DEFIRS.}
\l{}
\l{viatiko. I. Subsidio a persono por lua spensi dum lua voyajo.}
\l{\cel{3}- II. (religio katol.)\cel{2}Helpilo por acesar la vivo sequonta,}
\l{\cel{3}la vivo +nexta, eukaristio portata a persono extreme malada}
\l{\cel{3}por preparar lu pri morto bona. - DEFIS.}
\l{}
\l{vibrar. (netrans., de, pro). Agar ocilo duopla, movo alternanta}
\l{\cel{3}di la molekuli di korpo, qua transmisesas proximope a la medio.}
\l{\cel{4}DEFIS.}
\l{}
\l{vibriono. I.Ento mikroskopala, forme di filamento tenua. - II.}
\l{\cel{3}(biol.)\cel{2}(opoze a \textquotedbl{}bakterio\textquotedbl{}). Ento mikroskopala qua semblas}
\l{\cel{4}avancar per movo vermo-forma.}
}
